% Decision note (not thesis text): why the Riley sample-size criteria were taken out of the thesis.
% Written 2026-11-16 on request ("record the reason behind the changes in a .tex and compiled pdf").
% Sources of the numbers:
%   reports/Geometry feature budget cloglog GLMM.md (P_max table, own calculation of 2026-11-02, not
%     checked against pmsampsize; n = 1,076 participants, 452 events);
%   CGPS_data/saved.md (segment events, output of segment_k_budget.py, 4 November 2026);
%   root CLAUDE.md (the model: six risk factors, three geometric features, eight segment intercepts).
% The description of the Riley criteria is taken from the retired thesis section
%   (outdated/thesis_retired/technical_background/retired/01b_sample_size_2026-11-16.tex), whose
%   header says that Riley et al. 2020 was NOT opened in full. It is not verified against the paper.
% Everything in a "My reasoning" box is argument made for this project, not a sourced claim.
% All counts are from the 1,801 table, before image-quality exclusions and before the tracer cohort
%   is linked. They are preliminary.
\documentclass[11pt]{article}
\usepackage[a4paper,margin=25mm]{geometry}
\usepackage[T1]{fontenc}
\usepackage{lmodern}
\usepackage{amsmath,amssymb}
\usepackage{booktabs,array}
\usepackage[hidelinks]{hyperref}
\usepackage{parskip}
\usepackage{xcolor}
\usepackage{tcolorbox}
\newtcolorbox{decision}{colback=blue!4,colframe=blue!40!black,boxrule=0.4pt,left=4pt,right=4pt,top=3pt,bottom=3pt,title={Decision},fonttitle=\small\bfseries}
\newtcolorbox{mine}{colback=orange!6,colframe=orange!60!black,boxrule=0.4pt,left=4pt,right=4pt,top=3pt,bottom=3pt,title={My reasoning (not a sourced claim)},fonttitle=\small\bfseries}
\newtcolorbox{caution}{colback=red!4,colframe=red!55!black,boxrule=0.4pt,left=4pt,right=4pt,top=3pt,bottom=3pt,title={To be open about},fonttitle=\small\bfseries}

\title{Why the Riley sample-size criteria left the thesis}
\author{Decision note, 16 November 2026}
\date{}

\begin{document}
\maketitle

\begin{decision}
The Riley criteria are not used in the thesis, in the main text or in an appendix.
The section \emph{Sample Size for Prediction Models} of the Technical Background is retired.
The Methodology states the number of parameters of the model and the participants and events behind
them, and claims no threshold.
The decision was made on 16 November 2026, without the supervisor.
\end{decision}

\section{What the criteria are}

The retired section described them as follows.
Riley et al.\ (2020) give the sample size needed to develop a prediction model for a binary
outcome from three criteria.
The coefficients should need to be shrunk by no more than 10\,\%, so the expected shrinkage $S$ is at
least 0.9.
The Nagelkerke $R^2$ in the development data should be at most 0.05 above its optimism-adjusted value.
The overall proportion with the outcome should be estimated to within $\pm 0.05$.
Turned around, the first criterion gives the largest number of parameters a fixed sample supports,
\[
  P_{\max} = n\,(S-1)\,\ln\!\left(1 - \frac{R^2_{\mathrm{CS}}}{S}\right),
\]
where $R^2_{\mathrm{CS}}$ is the Cox--Snell $R^2$ the model is anticipated to reach, derived from an
anticipated C-statistic.

This description has not been verified against the paper.
The header of the retired section records that the paper could not be opened in full.

\section{The reasons}

\subsection*{1. The criteria protect predicted risks, and the main result is not a predicted risk}

Since 15 November 2026 the thesis is a prognostic factor study.
Its main result is one likelihood ratio test of three prespecified geometric terms, and the hazard
ratios of those terms.
All three criteria concern a model that will be used to predict: how extreme its predicted risks
are in new persons, how optimistic its apparent performance is, and how precisely it estimates the
average risk.

\begin{mine}
A test and hazard ratios depend on power and precision.
The number of events and the spread of each feature set these, and the confidence intervals in the
Results show them directly.
The covariates were fixed in advance and none is selected using the outcome, so the test does not
rest on the model being free of overfitting in Riley's sense.
\end{mine}

\subsection*{2. The secondary estimates are already cross-validated}

The change in C and the calibration are prediction quantities, and overfitting does matter for them.
They are estimated by cross-validation split by person.

\begin{mine}
Cross-validation measures the optimism of these estimates in the data.
A criterion that bounds the expected optimism in advance adds nothing once the optimism is measured.
The cross-validated calibration slope is the empirical counterpart of the shrinkage $S$ that the
first criterion tries to keep above 0.9.
\end{mine}

\subsection*{3. The criteria are derived for another setting}

The criteria assume independent persons with one binary outcome each.
The model is a segment-level complementary log-log mixed model: nine segments per participant, eight
segment intercepts and a random intercept per participant.
The earlier report found no published extension of the criteria to clustered data, so applying them
required a counting rule made for this project (which parameters count against the person-level
sample size).
The answer depended almost entirely on that rule and on the anticipated C, as the next section shows.

\subsection*{4. The criteria did not decide anything}

The three geometric features were chosen from the literature, and the six risk factors are the ones
SCORE2 uses.
The cohort size is fixed.
A ``parameter budget'' in the Methodology would therefore describe a decision process that did not
take place, as a check made after the model was already set.

\subsection*{5. The sources were not verified}

Riley et al.\ 2020 was not opened in full, and the paper behind the conversion from the C-statistic
(Riley, Van Calster and Collins 2021) is not in \texttt{refs.bib} and was not read.
The \texttt{pmsampsize} package was never run, because R is not installed.
The citation rules of the project do not allow a section to rest on sources in that state.
This reason alone would not justify the decision, since the papers could be obtained.
It is the cost of keeping a section that reasons 1 to 4 make unnecessary.

\section{What the criteria said about this model}

The table is copied from the report \emph{Geometry feature budget cloglog GLMM} (2 November 2026).
It is an own calculation of the first criterion at $S = 0.9$ for $n = 1{,}076$ participants and 452
events, and it was not checked against \texttt{pmsampsize}.

\begin{center}
\begin{tabular}{@{}lcc@{}}
\toprule
Anticipated C & $R^2_{\mathrm{CS}}$ & $P_{\max}$ \\
\midrule
0.65 & 0.068 & 8.4 \\
0.68 & 0.096 & 12.2 \\
0.70 & 0.118 & 15.1 \\
0.72 & 0.142 & 18.5 \\
0.75 & 0.181 & 24.2 \\
\bottomrule
\end{tabular}
\end{center}

The model has 9 covariates (six risk factors and three geometric features) and 8 segment intercepts,
17 fixed parameters in all.
If log follow-up, dominance and log segment length stay in the base, which is still open, it has 20.

\begin{center}
\begin{tabular}{@{}lccc@{}}
\toprule
What is counted & Parameters & Fits at C = 0.65 & Fits at C = 0.70 \\
\midrule
Covariates only & 9 & no (8.4) & yes (15.1) \\
Covariates and segment intercepts & 17 & no & no \\
All, with the three open base terms & 20 & no & no \\
\bottomrule
\end{tabular}
\end{center}

\begin{caution}
By the strict count the model is over the Riley limit even at an anticipated C of 0.70, and by the
most lenient count it is just over the limit at 0.65.
The line ``at the edge of the Riley criteria at C = 0.65'' in the project instructions was the lenient
count.
The decision was made knowing these numbers, so it sets aside a check that the model does not pass
on the strict count.
It is defensible only if reasons 1 to 4 hold on their own, and they do not depend on the outcome of
the check.
The decision is not a claim that the model is small enough by Riley's standard.
\end{caution}

The same report explains why the two counts differ.
The segment intercepts are estimated from where the events fall within a participant, about 970
segment events, and not from who has them.
Its simulation of the actual pooled model agreed with the count that leaves the segment intercepts out.
The project's own base model had an out-of-fold C of 0.687, between the two columns.
That figure comes from the person-level run, which is not the pooled model, so it is only an
indication.

\section{What the thesis says instead}

The section \emph{Parameter Budget and Feature Selection} of the Methodology is renamed
\emph{Model Size and Feature Selection}.
Its first subsection will state that the covariates were fixed before the geometric features were
related to the outcome, give the number of fixed parameters, and give the participants and events
behind them.
With the preliminary counts this is 452 participants who developed plaque in an analysed segment
among 1,076, and about 970 segment events.

\begin{center}
{\small\begin{tabular}{@{}lcc@{}}
\toprule
Parameters & Participants with the outcome, per parameter & Segment events, per parameter \\
\midrule
9 & 50 & 108 \\
17 & 27 & 57 \\
20 & 23 & 49 \\
\bottomrule
\end{tabular}}
\end{center}

The thesis gives the counts and no threshold such as ``10 events per parameter''.
A threshold would need a source that has been opened and that covers a clustered model.
The segment events are not independent, so the right-hand column overstates the information: the
earlier report estimated a design effect of about 3 at the within-person correlation it found most
plausible.

The precision of the hazard ratios is shown by their confidence intervals in the Results.
The bifurcation angle and the daughter asymmetry enter only the proximal LAD and proximal LCx rows,
so their intervals rest on fewer events than the interval for curvature.
This belongs in the Discussion.

\section{A second analysis that changes}

The secondary analysis of curvature shape planned the number of components $K$ as the smaller of the
number of components stable under bootstrap and the Riley limit of the segment.
The cap is removed, and $K$ is set by bootstrap stability alone.
The analysis is descriptive and restricted to the proximal and mid LAD, which had 275 and 201 events.
At an anticipated C of 0.65 the removed cap allowed 5 and 4 components there, and 10 and 8 at 0.70,
so it is unlikely to have been the binding limit.
If more components than that turn out to be stable, the events per parameter are stated next to $K$.

\section{What was changed}

{\small
\begin{center}
\begin{tabular}{@{}>{\raggedright\arraybackslash}p{0.44\linewidth}>{\raggedright\arraybackslash}p{0.50\linewidth}@{}}
\toprule
File & Change \\
\midrule
\path{content/technical/01b_sample_size.tex} & Moved to \path{outdated/thesis_retired/technical_background/retired/01b_sample_size_2026-11-16.tex} \\
\path{content/background.tex} & Input line removed \\
\path{content/technical/00_chapter_intro.tex} & Header comment updated \\
\path{content/methodology/04_model_size.tex} & Titles, notes and header: model size without Riley \\
\path{content/methodology/00_chapter_intro.tex} & Stage (6) reworded \\
\path{content/methodology/07_curvature_fpca.tex} & Header notes: $K$ by bootstrap stability alone \\
\path{content/discussion.tex} & Header comment: ``within the Riley budget'' removed \\
Root \path{CLAUDE.md} and the plan in \path{CGPS_statistics-/CLAUDE.md} & Decision recorded \\
\bottomrule
\end{tabular}
\end{center}
}

All thesis paths are under \texttt{Thesis\_coronary/}.
Three scripts still compute budgets and are left unchanged and unused:
\path{segment_k_budget.py}, \path{who_where_readiness.py} and \path{thesis_data.py}.
The entry \path{riley2020samplesize} stays in \path{refs.bib}, uncited.
To restore the section, move the retired file back and add its input line to
\texttt{content/background.tex}.

\section{Questions an examiner may ask}

\emph{How was the size of the model decided?}
The covariates were fixed before the geometric features were related to the outcome: the risk
factors are those of SCORE2 and the features are those the literature ties to low wall shear stress.
The size follows from that choice.

\emph{Is the model overfitted?}
The secondary performance estimates are cross-validated by person, and the cross-validated
calibration slope shows how much the predicted risks would need to be shrunk.
The thesis does not present the model as a tool for prediction.

\emph{Why no formal sample-size calculation?}
The cohort is fixed, so no calculation could change it.
The available criteria are for developing a prediction model from independent persons, and by the
strict count this model would not meet them, which the Discussion can state as a limit on how far
the secondary estimates can be trusted.

\section{Still open}

The counts must be replaced by those of the tracer cohort after the image-quality exclusions.
Whether log follow-up, dominance and log segment length stay in the base decides between 17 and 20
parameters.
Whether the Discussion states the strict Riley count as a limitation, as the last answer above
suggests, has not been decided.
It would bring one citation of Riley et al.\ 2020 back, and the paper would have to be read first.

\end{document}
